\section*{Bab \ref{ch:g10:soil-formation} --- Tanah, Sumber Daya yang Lambat Terbentuk}

\begin{solution}{exo:g10:soil-formation:1}
\omterm{def:g10:soil-formation:chemical-weathering}{Pelapukan kimia} adalah perubahan mineral-mineral batuan melalui reaksi dengan
air dan gas-gas yang dibawanya (karbon dioksida, oksigen): sebagian mineral
melarut, sebagian lain menjadi mineral baru yang stabil di permukaan.
\omterm{def:g10:soil-formation:hydrolysis}{Hidrolisis} adalah \omterm{def:g10:soil-formation:chemical-weathering}{pelapukan kimia} silikat oleh air dan ion-ion hidrogennya:
ion-ion logam masuk ke dalam larutan dan sisa mineral itu tersusun ulang
menjadi \omterm{def:g10:soil-formation:clay}{mineral lempung}.
\end{solution}

\begin{solution}{exo:g10:soil-formation:2}
\omterm{def:g10:soil-formation:weathering-profile}{Saprolit} adalah batuan yang terurai di tempatnya: lunak dan bersifat tanah,
\omterm{def:g10:soil-formation:weathering-profile}{saprolit} mempertahankan butir-butir, pita-pita, dan struktur \omterm{def:g10:soil-formation:weathering-profile}{batuan induknya}
serta beralih ke bawah, tanpa terputus, menjadi batuan segar. Tanah yang
terbentuk pada pasir sungai terletak di atas bahan yang dibawa dari tempat
lain: butir-butirnya membulat dan terpilah, dan tanah itu sama sekali tidak
memperlihatkan struktur batuan di bawahnya.
\end{solution}

\begin{solution}{exo:g10:soil-formation:3}
Feldspar menjadi lempung (kaolinit), sambil melepaskan kalium, natrium, atau
kalsium serta silika ke dalam air; mika hitam menjadi lempung dan oksida besi;
kuarsa tertinggal tanpa berubah, sebagai pasir.
\end{solution}

\begin{solution}{exo:g10:soil-formation:4}
Besi dalam mineral-mineral gelap, \ce{Fe^2+}, dioksidasi oleh oksigen yang
dibawa air menjadi oksida besi(III) yang tidak larut: goetit yang kuning
kecokelatan dan hematit yang merah, yang menyelimuti butir-butir dan mewarnai
seluruh bahan itu.
\end{solution}

\begin{solution}{exo:g10:soil-formation:5}
\qty{0.017}{mm} setahun dan \qty{1.537}{mm} setahun: $1.537/0.017 \approx 90$
kali lebih besar.
\end{solution}

\begin{solution}{exo:g10:soil-formation:6}
Kalsium: 1 dan 1. Karbon: $1 + 1 = 2$ dan 2. Oksigen: $3 + 2 + 1 = 6$ dan
$2 \times 3 = 6$. Hidrogen: 2 dan 2. Muatan: 0, dan $+2 - 2 = 0$.
\end{solution}

\begin{solution}{exo:g10:soil-formation:7}
Pelarutan kalsit hanya menghasilkan ion-ion terlarut, \ce{Ca^2+} dan
\ce{HCO3-}: semuanya terbawa pergi dan batu kapur murni tidak meninggalkan
sisa padat. \omterm{def:g10:soil-formation:hydrolysis}{Hidrolisis} feldspar menghasilkan padatan, kaolinit, yang tetap
tinggal; dan kuarsa dalam granit hampir tidak bereaksi sama sekali. Karena
itu granit meninggalkan lempung dan pasir.
\end{solution}

\begin{solution}{exo:g10:soil-formation:8}
Keping-keping lempung dan humus membawa muatan negatif yang menahan ion-ion
kalium positif \ce{K+} di permukaannya, sehingga air hujan yang melintas tidak
membawanya pergi. Akar melepaskan ion-ion hidrogen, yang menggantikan kalium
pada keping-keping itu; kalium yang terlepas ke dalam larutan tanah diserap oleh
akar.
\end{solution}

\begin{solution}{exo:g10:soil-formation:9}
\omterm{def:g10:soil-formation:chemical-weathering}{Pelapukan kimia} memerlukan air, yang berlimpah di iklim basah dan langka di
iklim kering (dan membeku di iklim dingin); dan reaksi-reaksinya berlangsung
lebih cepat bila hangat. Hanya di bawah panas dan air yang berlimpah, dalam
waktu yang lama, reaksi-reaksi itu berlangsung cukup jauh untuk menyingkirkan
hampir semua yang dapat larut dan meninggalkan oksida besi dan aluminium.
\end{solution}

\begin{solution}{exo:g10:soil-formation:10}
$10/0.017 \approx 590$ tahun untuk \qty{1}{cm}; $300/0.017 \approx 17\,600$
tahun untuk \qty{30}{cm}; $1000/0.017 \approx 59\,000$ tahun untuk \qty{1}{m}.
\end{solution}

\begin{solution}{exo:g10:soil-formation:11}
Volume: $10\,000 \times 0.0015 = \qty{15}{m^3}$; massa: $15 \times 1200 =
\qty{18000}{kg}$, \num{18}~ton per hektare per tahun.
\end{solution}

\begin{solution}{exo:g10:soil-formation:12}
$46\,000/0.017 \approx 2.7$ juta tahun. \omterm{def:g10:soil-formation:weathering-profile}{Saprolit} setebal itu memerlukan bentang
alam yang telah stabil, tanpa selimut lapuknya terkikis habis, selama jutaan
tahun.
\end{solution}

\begin{solution}{exo:g10:soil-formation:13}
Goetit: \ce{4 Fe^2+ + O2 + 6 H2O -> 4 FeO(OH) + 8 H+}: besi 4 dan 4,
oksigen $2 + 6 = 8$ dan $4 \times 2 = 8$, hidrogen 12 dan $4 + 8 = 12$, muatan
$+8$ dan $+8$. Hematit: \ce{4 Fe^2+ + O2 + 4 H2O -> 2 Fe2O3 + 8 H+}: besi 4
dan 4, oksigen $2 + 4 = 6$ dan 6, hidrogen 8 dan 8, muatan $+8$ dan $+8$.
\end{solution}

\begin{solution}{exo:g10:soil-formation:14}
Kehilangan bersih $0.082 - 0.017 = \qty{0.065}{mm}$ setahun; $300/0.065 \approx 4600$
tahun. Di bawah pembajakan konvensional, $300/1.52 \approx 200$ tahun: lebih
dari dua puluh kali lebih singkat.
\end{solution}

\begin{solution}{exo:g10:soil-formation:15}
Satu ton itu mengandung \qty{300}{kg} ortoklas. Setiap kilogram menghasilkan
\qty{464}{g} kaolinit: $300 \times 0.464 \approx \qty{139}{kg}$. Kalium:
$300 \times 0.1405 \approx \qty{42}{kg}$.
\end{solution}

\begin{solution}{pb:g10:soil-formation:1}
\textbf{1.} $4 \times 0.3048 \approx \qty{1.22}{m}$.

\textbf{2.} $0.5 \times 30.48 \approx \qty{15}{cm}$.

\textbf{3.} Humus, bahan organik gelap yang dibuat oleh pengurai dari
tumbuhan mati.

\textbf{4.} Oksida besi, goetit dan hematit, yang terbentuk oleh oksidasi
besi(II) dalam mineral-mineral bahan induk.

\textbf{5.} Tidak. Di bawahnya, profil itu akan beralih menjadi batuan lapuk,
yaitu \omterm{def:g10:soil-formation:weathering-profile}{saprolit} (atau bahan lapuk lainnya), lalu menjadi \omterm{def:g10:soil-formation:weathering-profile}{batuan induk} yang
segar.

\textbf{6.} \ce{2 KAlSi3O8 + 2 H+ + 9 H2O -> Al2Si2O5(OH)4 + 2 K+ + 4 H4SiO4}.

\textbf{7.} $1000/278.33 \approx \qty{3.59}{mol}$.

\textbf{8.} $3.59/2 \approx \qty{1.80}{mol}$ kaolinit, $1.80 \times 258.16
\approx \qty{464}{g}$.

\textbf{9.} Ion-ion kalium dan silika terlarut: keduanya dibawa air ke
sungai-sungai, dan oleh sungai ke laut.

\textbf{10.} Kuarsa hampir tidak bereaksi: butir-butir kuarsa tertinggal sebagai pasir.

\textbf{11.} $1219/0.017 \approx 72\,000$ tahun.

\textbf{12.} $1.537 - 0.017 = \qty{1.52}{mm}$ setahun.

\textbf{13.} $1219/1.52 \approx 800$ tahun.

\textbf{14.} $72\,000/800 = 90$: profil itu hilang sembilan puluh kali lebih
cepat daripada saat terbentuk.

\textbf{15.} Tidak: $0.013 < 0.017$, erosi sedikit lebih lambat daripada
produksi dan tanah itu bertahan.

\textbf{16.} $T = S/(E - P)$.

\textbf{17.} $300/1.52 \approx 200$ tahun.

\textbf{18.} $300/(0.082 - 0.017) = 300/0.065 \approx 4600$ tahun.

\textbf{19.} Tanah yang hilang dalam beberapa abad memerlukan puluhan ribu
tahun untuk terbentuk kembali: pada skala hidup manusia dan masyarakat, apa
yang hilang tidak tergantikan.

\textbf{20.} $1000/1.52 \approx 660$ tahun: \textbf{satu meter tanah hilang
dalam sekitar 660 tahun, beberapa abad, di bawah pembajakan konvensional}.
\end{solution}
